\begin{center}
  \textsc{Abstract}
\end{center}
%
\noindent
%
Clinical natural language processing relies on large volumes of annotated text, yet access to real patient records remains severely restricted by privacy regulations such as GDPR and HIPAA.
This thesis investigates synthetic data generation as a privacy-preserving alternative for training clinical NLP systems.

We propose a facsimile generation approach in which synthetic clinical documents are produced by a large language model guided solely by medical concept keywords extracted from real records via UMLS, ensuring that no patient text crosses the privacy boundary.
The generator is iteratively refined through Direct Preference Optimisation (DPO) using a private contrastive scorer that compares synthetic and real documents, returning only numerical scores to the generation side.
Evaluated on MIMIC-III discharge summaries, synthetic data achieves ICD-9 classification performance comparable to or exceeding that of real data, particularly on complex label spaces where oversampling compensates for the quality gap.
We further demonstrate the generalisability of the approach to French clinical reports from AP-HP, where DPO alignment yields substantial gains in ICD-10 coding performance, confirming the language-agnostic nature of the methodology.

In parallel, we explore LLM-based weak annotation as a complementary strategy for low-resource clinical NER.
Through knowledge distillation from a large language model into smaller encoder models on the multilingual E3C corpus, we show that hybrid supervision combining LLM-generated and dictionary-based annotations achieves the best performance across five European languages, with particular benefits for low-resource settings.

Finally, we conduct a thorough privacy analysis revealing that DPO alignment, while improving generation quality, amplifies re-identification risk.
We propose keyword perturbation as a mitigation strategy that substantially reduces linkage attack success at negligible utility cost for conversational tasks.
However, discrimination attacks remain partially effective even after mitigation, underscoring that no zero-risk solution exists.

This work demonstrates that keyword-guided synthetic generation with iterative alignment offers a viable path toward privacy-preserving clinical NLP, while honestly acknowledging the residual risks and the fundamental tension between data utility and patient privacy.
